\begin{mainquestion}{capacitornetwork}{Capacitor Network at Steady State}

\questionpoints Consider the capacitor network shown below.

Initially, the voltage between nodes $A$ and $B$ is zero, and the voltage
across every capacitor is $0\text{ V}$.

At time $t=0$, a $\SI{1}{\milli\ampere}$ current source is turned on between
nodes $A$ and $B$. The current source remains on until the voltage between
$A$ and $B$ reaches $\SI{15}{\volt}$, at which point it is turned off.

\begin{center}
\begin{circuitikz}[american, scale=0.9, transform shape]
    % Top and Bottom rails
    \draw (-6,3) node[left]{$A$} -- (6,3);
    \draw (-6,-3) node[left]{$B$} -- (6,-3);

    % Leg 1: C1 and C6 in series
    \draw (-5,3) to[C, l=$1\,\mu\text{F}$] (-5,0)
                to[C, l=$6\,\mu\text{F}$] (-5,-3);

    % Leg 2: C7 alone
    \draw (-2.5,3) -- (-2.5,0)
                   to[C, l=$7\,\mu\text{F}$] (-2.5,-3);

    % Leg 3: C2, C4, C8 in series
    \draw (0,3) to[C, l=$2\,\mu\text{F}$] (0,1)
               to[C, l=$4\,\mu\text{F}$] (0,-1)
               to[C, l=$8\,\mu\text{F}$] (0,-3);

    % Leg 4: C3 alone
    \draw (2.5,3) -- (2.5,0)
                   to[C, l=$3\,\mu\text{F}$] (2.5,-3);

    % Leg 5: C5 alone
    \draw (5,3) -- (5,0)
                to[C, l=$5\,\mu\text{F}$] (5,-3);

    % Connection dots
    \foreach \x in {-5,-2.5,0,2.5,5} {
        \node[circ] at (\x,3) {};
        \node[circ] at (\x,-3) {};
    }
\end{circuitikz}
\end{center}

% =========================================================
\begin{partquestion}{sourceduration}
{How long does the current source remain on?}

\begin{answer}[1.2in]
The total effective capacitance between nodes $A$ and $B$ is found by first
computing the effective capacitance of each parallel leg.

\medskip
\textbf{Leg 1: $1\,\mu\text{F}$ and $6\,\mu\text{F}$ in series}
\[
C_{eq,1} = \left(\frac{1}{1} + \frac{1}{6}\right)^{-1}
= \frac{6}{7}\,\mu\text{F}
\]

\textbf{Leg 2: $7\,\mu\text{F}$}
\[
C_{eq,2} = 7\,\mu\text{F}
\]

\textbf{Leg 3: $2\,\mu\text{F}$, $4\,\mu\text{F}$, and $8\,\mu\text{F}$ in series}
\[
C_{eq,3} =
\left(\frac{1}{2} + \frac{1}{4} + \frac{1}{8}\right)^{-1}
= \frac{8}{7}\,\mu\text{F}
\]

\textbf{Leg 4: $3\,\mu\text{F}$}
\[
C_{eq,4} = 3\,\mu\text{F}
\]

\textbf{Leg 5: $5\,\mu\text{F}$}
\[
C_{eq,5} = 5\,\mu\text{F}
\]

\medskip
Since the five legs are in parallel, the total capacitance is the sum:
\[
C_\text{total}
= \frac{6}{7} + 7 + \frac{8}{7} + 3 + 5
= \boxed{17\,\mu\text{F}}.
\]

The voltage across the network increases according to
\[
I = C_\text{total}\frac{dV}{dt}.
\]

Solving for the time required to reach $\SI{15}{\volt}$:
\[
t = \frac{C_\text{total} V}{I}
= \frac{(17\,\mu\text{F})(15\,\text{V})}{1\,\text{mA}}
= \boxed{\SI{0.255}{\second}}.
\]
\end{answer}

\end{partquestion}

% =========================================================
\begin{partquestion}{finalvoltages}
{Determine the final voltage across each capacitor.

\emph{Hint: Capacitors connected in series carry the same current.
The current through a given series stack depends on the effective
capacitance of that stack.}}

\begin{answer}
Each of the five parallel legs has a total voltage of $15\text{V}$ across it.

\medskip
\textbf{Single-capacitor legs:}

Since each of these legs contains only one capacitor:
\[
V_{C7} = V_{C3} = V_{C5} = 15\text{V}
\]

\medskip
\textbf{First leg: $C_1$ and $C_6$ in series}

For series capacitors, the charge is equal:
\[
Q_1 = Q_6 \quad \Rightarrow \quad C_1 V_{C1} = C_6 V_{C6}
\]

Also:
\[
V_{C1} + V_{C6} = 15
\]

Solving:
\[
V_{C1} = \frac{15 C_6}{C_1 + C_6}
       = \frac{15 \cdot 6}{1 + 6}
       = \boxed{\frac{90}{7}\text{ V}}
\]

\[
V_{C6} = 15 - \frac{90}{7}
       = \boxed{\frac{15}{7}\text{ V}}
\]

\medskip
\textbf{Third leg: $C_2$, $C_4$, and $C_8$ in series}

Equal charge condition:
\[
C_2 V_{C2} = C_4 V_{C4} = C_8 V_{C8}
\]

And:
\[
V_{C2} + V_{C4} + V_{C8} = 15
\]

This gives:
\[
V_{C2} = \frac{15 C_4 C_8}{C_2 C_4 + C_2 C_8 + C_4 C_8}
\]

Substitute values:
\[
V_{C2} = \frac{15(4)(8)}{(2)(4) + (2)(8) + (4)(8)}
       = \frac{480}{56}
       = \boxed{\frac{60}{7}\text{ V}}
\]

Then:
\[
V_{C8} = \frac{C_2}{C_8} V_{C2}
       = \frac{2}{8} \cdot \frac{60}{7}
       = \boxed{\frac{15}{7}\text{ V}}
\]

\[
V_{C4} = \frac{C_2}{C_4} V_{C2}
       = \frac{2}{4} \cdot \frac{60}{7}
       = \boxed{\frac{30}{7}\text{ V}}
\]

\medskip
\textbf{Final answers:}

\[
\begin{aligned}
V_{C1} &= \frac{90}{7}\text{ V}, & V_{C6} &= \frac{15}{7}\text{ V} \\
V_{C2} &= \frac{60}{7}\text{ V}, & V_{C4} &= \frac{30}{7}\text{ V}, & V_{C8} &= \frac{15}{7}\text{ V} \\
V_{C3} &= 15\text{ V}, & V_{C5} &= 15\text{ V}, & V_{C7} &= 15\text{ V}
\end{aligned}
\]
\end{answer}

\end{partquestion}

\end{mainquestion}
